\documentclass[11pt]{article}
\usepackage[margin=1.1in]{geometry}
\usepackage{times}
\usepackage[T1]{fontenc}
\usepackage{setspace}
\onehalfspacing
\usepackage{parskip}
\pagenumbering{gobble}

\title{The Last Human Job}
\author{Flyxion}
\date{September 2026}

\begin{document}
\maketitle

Every occupation with a describable output was automated decades ago, cleanly, without much resistance, because each one went individually and the ones still employed always had a good reason to believe their own case was different until the week it wasn't. What remains, the single human role left on the org chart, is a chair in a small room facing a console with one button, unlabeled except for a strip of engraved plastic that reads AUTHORIZE, and a person, hired through a competitive process, whose entire function is to press it when a light above it turns amber.

Nobody currently employed knows what the button authorizes. The systems downstream of it were redesigned so many revisions ago that the chain connecting the press to whatever it triggers has been rerouted, abstracted, and optimized past the point where any single engineer holds the whole picture, and the honest answer, if the person in the chair asked, which the job description discourages, is that no one does either. The button is pressed several times a shift, always after the amber light, always followed by a green confirmation, and the entire apparatus downstream continues operating exactly as it would if the chair were empty, which it very nearly was during a staffing shortfall two years back, a fact quietly excluded from the annual report.

The position is described, in recruiting material and in testimony to the regulatory bodies that require its existence, as evidence of meaningful human oversight retained at the center of an otherwise fully automated system, and in a narrow sense this is true: a human is there, a human presses it, a human could, in principle, refuse. No one has ever refused. No one has been trained on what refusing would do, or given information sufficient to decide whether it should be done, and the job's entire value to the organization that funds it rests on this specific arrangement never changing, because a chair that could meaningfully say no is a much more expensive thing to staff than a chair that has simply never been asked to.

\end{document}
